% Slide 4 — Consequências operacionais
\begin{frame}{Consequências de uma política única}
  \begin{columns}[T,onlytextwidth]
    \column{0.33\textwidth}
      \begin{block}{Custo}
        \small Excesso de estoque parado, ocupando espaço e capital.
      \end{block}
    \column{0.33\textwidth}
      \begin{block}{Ruptura}
        \small Falta de produto na prateleira, venda perdida, cliente frustrado.
      \end{block}
    \column{0.33\textwidth}
      \begin{block}{Efeito chicote}
        \small Pedidos das lojas amplificam a variabilidade observada pelo
        armazém regional.
      \end{block}
  \end{columns}
  \vfill
  \begin{center}
    \small Custo, disponibilidade e estabilidade dos pedidos são objetivos
    \highlight{em conflito} --- e o equilíbrio adequado depende do
    \highlight{perfil da série}.
  \end{center}
  \note{
    Quando uma política única é aplicada indiscriminadamente, três
    consequências aparecem. Primeiro, excesso de custo de manutenção de
    estoque. Segundo, ruptura: falta de produto e venda perdida. Terceiro,
    e menos óbvio, o efeito chicote: pequenas variações de demanda na loja
    podem virar oscilações muito maiores nos pedidos enviados ao armazém,
    afetando a cadeia inteira. Esses três objetivos -- custo, serviço e
    estabilidade -- estão em conflito, e o ponto de equilíbrio ideal não é
    o mesmo para todas as séries. Isso é o que vamos quantificar mais
    adiante nos experimentos.
  }
\end{frame}
